% !PRINTED 789
\weberpage{789}

\noindent
Gruppen, \gesperrt{Abel}'sche I, 476, 536; II, 6.\\
--- endliche II, 4.\\
--- von Permutationen I, 475.\\
--- von Substitutionen I, 472.\\
--- vom Grade $p^{\alpha}$ II, 127.\\
--- vom Grade $p\,q$ II, 141.\\
--- linearer Substitutionen II, 152.\\
--- orthogon. Substitutionen II, 184.\\
--- linearer gebrochener Substitutionen II, 189, 195.\\
Gruppencharaktere II, 44.\\
Gruppentafel II, 115.

\medskip
\noindent
\textbf{H}albmetacyklische Gruppen I, 616.\\
Hauptaxen einer ternären linearen Substitution II, 441.\\
--- der Gruppe $G_{168}$ II, 444.\\
Hauptgleichung fünften Grades I, 175, 230.\\
Hauptreihe II, 24.\\
Hauptunterdeterminanten I, 256.\\
\gesperrt{Hermite}'s Lösung des \gesperrt{Sturm}'schen Problems I, 280.\\
\gesperrt{Hesse}'sche Curve II, 330.\\
--- Determinante I, 186.\\
\gesperrt{Hesse-Cayley}'sche Bezeichnung der Doppeltangenten II, 369.\\
\gesperrt{Hilbert}'scher Satz II, 165.\\
Homogene Functionen I, 56.

\medskip
\noindent
\textbf{I}deale II, 547.\\
Ideale Zahlen nach \gesperrt{Kummer} II, 551.\\
Identische Gruppe I, 496.\\
--- Substitution II, 152.\\
Ikosa\"edergleichung II, 232, 418.\\
Ikosa\"edergruppe II, 220, 279.\\
Ikosa\"ederinvarianten II, 232.\\
Ikosa\"ederresolventen II, 422, 426.\\
Imaginäre Zahlen I, 17.\\
Imaginäre Form der linearen Congruenzgruppe II, 266.\\
Imaginärer Congruenzkörper zweiten Grades II, 259.\\
Imaginäre quadratische Irrationalzahlen I, 379.\\
Imprimitive Gruppen I, 483, 516.\\
--- Körper I, 460, 483.\\
Index einer Invariante II, 163.\\
--- eines Theilers einer Gruppe I, 502; II, 8.

Indexmoduln II, 61.\\
Indexreihe einer Gruppe II, 17.\\
Indices I, 431.\\
--- der Elemente einer Gruppe II, 44.\\
--- nach einer ungeraden Primzahlpotenz II, 54.\\
--- nach einer Potenz von 2 II, 58.\\
--- nach einem zusammengesetzten Modul II, 60.\\
Inflexionspunkte II, 326.\\
--- einer Curve dritter Ordnung II, 331.\\
Interpolation I, 44, 98, 331.\\
Intransitive Gruppen I, 482.\\
--- Normaltheiler I, 520.\\
Invarianten I, 186.\\
--- der binären cubischen Form I, 193.\\
--- --- --- biquadratischen Form I, 200, 206.\\
--- einer \gesperrt{Abel}'schen Gruppe II, 42.\\
--- endlicher Gruppen linearer Substitutionen II, 161.\\
--- absolute und relative II, 162, 177.\\
--- im weiteren Sinne II, 179.\\
--- der ternären Formen II, 328.\\
--- der Curven dritter Ordnung II, 337.\\
--- der Gruppe $G_{168}$ II, 453, 461.\\
Invariantencurven II, 454.\\
Inverse Substitution II, 155.\\
Irrationale Verhältnisse I, 10.\\
--- Zahlen I, 12, 370.\\
Irreducibilität I, 453.\\
--- reiner Gleichungen I, 607.\\
--- der Kreistheilungsgleichung I, 417; II, 768.\\
Isomorphe Gruppen I, 477; II, 6.\\
Isomorphismus (mehrstufiger) II, 15.

\medskip
\noindent
\gesperrt{\textbf{J}acobi}'s Abschätzung der Zahl der Wurzeln zwischen zwei Grenzen I, 311.\\
\gesperrt{Jordan}'scher Satz über zusammengesetzte Gruppen II, 18.

\medskip
\noindent
\textbf{K}ettenbrüche I, 358.\\
--- für äquivalente Zahlen I, 374.\\
Ketten von Hauptunterdeterminanten I, 257.\\
--- \gesperrt{Sturm}'sche I, 272.
